\documentclass[11pt,a4paper]{article}

\usepackage{../shared/cathedral-preamble}

\title{%
  The Cathedral and Society:\\
  Legal, Policy, and Educational Perspectives\\
  on Machine-Verified Mathematics%
}
\author{Jason Robert Gochanour}
\date{May 2026}

\begin{document}
\maketitle

\begin{abstract}
We examine the legal, policy, and educational implications of the
Cathedral project. We address intellectual property questions
(can a machine-verified proof be patented? who owns AI-assisted
mathematics?), governance implications (how should formally
verified results be treated by standards bodies?), and educational
applications (how can the Cathedral be used to teach analytic number
theory, formal verification, and mathematical research methodology?).
\end{abstract}

\tableofcontents

% ================================================================
\section{Introduction}
% ================================================================

The Cathedral is a mathematical object, but it exists in a legal,
political, and educational context. A formally verified reduction
of the Riemann Hypothesis raises questions that pure mathematics
does not usually encounter.

% ================================================================
\section{Intellectual Property}
% ================================================================

\subsection{Copyright}

Mathematical theorems cannot be copyrighted (Baker v. Selden, 1879).
However, the \emph{expression} of a proof---the specific Lean code,
the paper text, the documentation---is copyrightable. The Cathedral's
code is released under an open-source license.

\subsection{Patents}

Mathematical methods are generally not patentable (Alice Corp. v.
CLS Bank, 2014). However, specific \emph{applications} of the
Cathedral's mathematical infrastructure---such as the M\"obius
Harmonic Analyzer or the arithmetic fingerprinting protocol---could
potentially be patented as methods or devices.

The author has chosen not to pursue patents, consistent with the
tradition of mathematical openness.

\subsection{AI Authorship}

The Cathedral was built with significant AI contributions. Under
current U.S. law (Thaler v. Perlmutter, 2023), AI systems cannot
be listed as authors or inventors. The human author retains sole
authorship, with AI contributions acknowledged.

This legal framework may evolve. The Cathedral's detailed records
of human-AI collaboration (COMM-LINKs, exploration logs) provide
a useful case study for future legal analysis.

% ================================================================
\section{Policy Implications}
% ================================================================

\subsection{Standards for Formal Verification}

Should formally verified mathematical results carry more weight
than traditional peer-reviewed proofs? We argue yes, with caveats:

\begin{enumerate}
  \item \textbf{Compiler trust}: The verification is only as
    trustworthy as the Lean kernel (5{,}000 LOC).
  \item \textbf{Specification trust}: The formalization must
    faithfully represent the mathematical intent.
  \item \textbf{Axiom transparency}: All axioms must be
    explicitly documented.
\end{enumerate}

The Cathedral satisfies all three criteria.

\subsection{Dual-Use Governance}

The Cathedral's dual-use paper (\emph{The Dark Mirror}) identifies
7 risk categories where the mathematical infrastructure could be
weaponized. This raises governance questions:
\begin{itemize}
  \item Should verified instability certificates be classified?
  \item Who is responsible for misuse of openly published
    mathematical results?
  \item How should the tension between mathematical openness and
    security be managed?
\end{itemize}

We follow the tradition of responsible disclosure: publishing
the mathematics with explicit threat analysis and defense strategies.

\subsection{Funding and Incentives}

Formal verification is expensive (45 days of intensive work for
the Cathedral). Current academic incentive structures reward
traditional papers, not machine-verified proofs. Policy changes
that could help:
\begin{itemize}
  \item Funding agencies accepting formal proofs as deliverables.
  \item Journals accepting Lean code as supplementary material.
  \item Tenure committees recognizing formalization contributions.
\end{itemize}

% ================================================================
\section{Educational Applications}
% ================================================================

\subsection{Teaching Analytic Number Theory}

The Cathedral provides a structured curriculum:
\begin{enumerate}
  \item \textbf{Week 1--2}: The Nyman--Beurling criterion
    (definitions, Gram matrix, positive definiteness).
  \item \textbf{Week 3--4}: The converse direction
    (Mellin transforms, rank-1 factorization).
  \item \textbf{Week 5--6}: The forward direction
    (Parseval bridge, Hardy--Littlewood MVT).
  \item \textbf{Week 7--8}: Numerical experiments
    (Rust/MPFR, Vasyunin formula, eigenvalue computation).
\end{enumerate}

Each topic has both a mathematical treatment (the companion papers)
and a machine-verified formalization (the Lean code).

\subsection{Teaching Formal Verification}

The Cathedral demonstrates practical ITP skills:
\begin{itemize}
  \item Working with Mathlib APIs (interval integrals, matrices).
  \item Managing axioms in evolving codebases.
  \item The layered module architecture pattern.
  \item Debugging type-class inference failures.
\end{itemize}

\subsection{Teaching Research Methodology}

The Cathedral's archive of 35 explorations, including abandoned approaches,
5 formalization traps, and 20+ axiom graduations provides a realistic
picture of mathematical research---including the failures that are
usually invisible in published papers.

% ================================================================
\section{Accessibility}
% ================================================================

The Cathedral is designed for accessibility:
\begin{itemize}
  \item All code is open source.
  \item All papers are freely available.
  \item The Lean code compiles with standard tooling.
  \item The Rust experiments build with \texttt{cargo build}.
  \item No proprietary software is required at any step.
\end{itemize}

% ================================================================
\section{Conclusion}
% ================================================================

The Cathedral exists at the intersection of mathematics, computer
science, law, and education. Its open-source release, transparent
axiom documentation, and dual-use risk assessment provide a model
for responsible publication of formally verified frontier mathematics.

\begin{thebibliography}{99}
\bibitem{cathedral}
J.R.~Gochanour, \emph{The Cathedral}, 2026.
\bibitem{alice2014}
Alice Corp. v. CLS Bank International, 573 U.S. 208 (2014).
\bibitem{thaler2023}
Thaler v. Perlmutter, No. 1:22-cv-01564 (D.D.C. 2023).
\end{thebibliography}

\end{document}
